%% ma: L11-B17-0061
%% lop: 11
%% chuong: 5
%% bai: 17
%% dang: 11.17.1
%% loai: DS
%% muc: [TH, TH, NB, TH]
%% dap_an: "ĐĐSĐ"
%% nguon: "tai-lieu-hoc-tap-mon-toan-lop-11-theo-chuong-trinh-giao-duc-pho-thong-2018.pdf (Lớp Toán Thầy Duy, Phạm Lê Duy, THPT Chu Văn An), Chương 3 Bài 3, Phần 2, Câu 1 (tr.414)"
%% kien_thuc: "Lớp 11 Bài 17 (hàm số liên tục)"
%% trang_thai: cho-duyet
%% ly_do: "Dạng mới đề xuất (11.17.x), chờ giáo viên duyệt dạng; tài liệu gốc không có lời giải mẫu, đáp án do kho tự giải và kiểm bằng sympy."
\begin{ex}
Cho hàm số $f(x)=\begin{cases}-\dfrac x2&\text{khi }x\le1\\\dfrac{x^2-3x+2}{x^2-1}&\text{khi }x>1\end{cases}$ và $g(x)=x^2-3x+1$.
\choiceTF
{\True Hàm số $f(x)$ liên tục tại điểm $x_0=1$.}
{\True Hàm số $g(x)$ liên tục tại điểm $x_0=1$.}
{Giới hạn $\lim\limits_{x\to1^+}f(x)=\dfrac12$.}
{\True Hàm số $y=f(x)+g(x)$ liên tục tại điểm $x_0=1$.}
\loigiai{\sloppy $f(1)=-\dfrac12$, $\lim\limits_{x\to1^-}f(x)=-\dfrac12$, $\lim\limits_{x\to1^+}\dfrac{(x-1)(x-2)}{(x-1)(x+1)}=\lim\dfrac{x-2}{x+1}=-\dfrac12$.
a) $f$ liên tục tại $1$. Đúng.
b) $g$ là đa thức nên liên tục. Đúng.
c) $\lim\limits_{x\to1^+}f(x)=-\dfrac12\ne\dfrac12$. Sai.
d) Tổng của hai hàm liên tục tại $1$ liên tục tại $1$. Đúng.}
\end{ex}
